\def\EditionMode{1}
\def\EditionLabel{Local math-prose proof}
\def\EditionStatus{Proposed wrappers; canonical source unchanged}
\documentclass[10pt,oneside]{book}
\usepackage{fontspec}
\setmainfont{Latin Modern Roman}
\setsansfont{Latin Modern Sans}
\usepackage{amsmath,amssymb,mathtools}
\usepackage{unicode-math}
\setmathfont{Latin Modern Math}
\usepackage{xcolor,tikz-cd,lua-ul,marginnote,pdflscape}
\usetikzlibrary{arrows.meta,decorations.pathmorphing,calc,shapes.geometric}
\tikzcdset{ampersand replacement=\&}
\tikzcdset{every label/.append style={text=black}}
\usepackage[a4paper,left=18mm,right=45mm,top=18mm,bottom=22mm,marginparwidth=28mm,marginparsep=7mm]{geometry}
\usepackage[unicode,colorlinks=true,urlcolor=black,linkcolor=black]{hyperref}
\usepackage{bookmark}
\setlength{\parindent}{0pt}
\setlength{\parskip}{4pt plus 1pt minus 1pt}
\setlength{\emergencystretch}{1.5em}
\raggedbottom
% Include final-line descenders inside the declared text-height rectangle.
\maxdepth=0pt
\definecolor{PSHand}{RGB}{25,58,132}
\definecolor{PSDelete}{RGB}{146,37,37}
\definecolor{PSTypedOld}{RGB}{95,95,95}
\ifnum\EditionMode=2\colorlet{blue}{black}\else\colorlet{blue}{PSHand}\fi
\RenewDocumentCommand{\textcircled}{m}{\tikz[baseline=(PSCircle.base)]\node[draw,circle,inner sep=.5pt,line width=.4pt](PSCircle){#1};}
\newif\ifPSInApparatus
\newcounter{psnote}
\newcounter{psvariant}
\newcounter{pslinkmark}
\newif\ifPSInLink
\newsavebox{\PSNoteBox}
\newdimen\PSNoteLimit
{\footnotesize\global\PSNoteLimit=3\baselineskip}
\newcommand{\CurrentScanURL}{}
\newcommand{\CurrentKey}{}
\newcommand{\CurrentScan}{}
\newcommand{\pssource}[3]{\psflushmargins\clearpage\gdef\CurrentKey{#1}\gdef\CurrentScan{#2}\gdef\CurrentScanURL{#3}\setcounter{psvariant}{0}\typeout{PS-SOURCE: #1}}
\newcommand{\psdivision}[2]{\pdfbookmark[0]{#1}{#2}\par\noindent{\large\bfseries #1\par}\medskip}
\newcommand{\pslinkmarker}[2]{\stepcounter{pslinkmark}\latelua{texio.write_nl('log','PS-LINK-SHIP|#1|\CurrentKey|\arabic{pslinkmark}|#2')}}
\newcommand{\pslink}[2]{\ifPSInLink\pslinkmarker{#1}{nested}#2\else\begingroup\PSInLinktrue\href{\CurrentScanURL}{\pslinkmarker{#1}{outer}#2}\endgroup\fi}
\newcommand{\scanpage}[1]{\ifnum\EditionMode=2\leavevmode\marginpar{\scriptsize\sffamily\pslink{scanpage}{scan #1}}\else\par\noindent{\small\sffamily\pslink{scanpage}{[scan #1]}}\par\medskip\fi}
\newcommand{\mspage}[1]{\ifnum\EditionMode=2\leavevmode\marginpar{\scriptsize\sffamily p.\,#1}\else{\small\sffamily[p.\,#1]}\fi}
\newcommand{\hand}[1]{{\ifnum\EditionMode=2\color{black}\else\color{PSHand}\fi #1}}
\let\over\relax
\NewDocumentCommand{\over}{O{} m m}{{\textcolor{PSTypedOld}{#2}\ensuremath{^{\!\nearrow}}\if\relax\detokenize{#1}\relax\hand{#3}\else #3\fi}}
\NewDocumentCommand{\inserted}{O{} m}{{\if\relax\detokenize{#1}\relax\hand{\ensuremath{\langle}#2\ensuremath{\rangle}}\else\ensuremath{\langle}#2\ensuremath{\rangle}\fi}}
\newcommand{\deleted}[1]{{\color{PSDelete}\ifmmode\text{[del.: }#1\text{]}\else[\textsf{del.:} #1]\fi}}
\newcommand{\xout}[1]{{\color{PSTypedOld}\ifmmode\text{[x: }#1\text{]}\else[\textsf{x:} #1]\fi}}
\let\PSMathUnderline\underline
\RenewDocumentCommand{\underline}{+m}{\ifmmode\PSMathUnderline{#1}\else\underLine{#1}\fi}
\NewDocumentCommand{\uline}{+m}{\underline{#1}}
\NewDocumentCommand{\uuline}{+m}{\ifmmode\PSMathUnderline{\PSMathUnderline{#1}}\else\underLine{\underLine{#1}}\fi}
\newunderlinetype\PSWaveLine{\leaders\hbox{\lower.7ex\hbox{\tikz[x=1pt,y=1pt,baseline=0pt]{\draw[line width=.3pt] (0,0) .. controls (1,1) and (2,1) .. (3,0) .. controls (4,-1) and (5,-1) .. (6,0);}}}}
\NewDocumentCommand{\uwave}{+m}{{\PSWaveLine #1}}
\newcommand{\unsure}[1]{\pslink{unsure}{[#1\ifmmode\,?\else\,?\fi]}}
\newcommand{\illegible}[1]{\ifmmode\text{\pslink{illegible}{[illeg.\ #1]}}\else\pslink{illegible}{[illeg.\ #1]}\fi}
\NewDocumentCommand{\diagram}{+m}{\ifmmode\vcenter{\hbox{#1}}\else\par\medskip{\centering #1\par}\medskip\psflushmargins\fi}
\newenvironment{diagramblock}{\par\begin{center}}{\end{center}\par}
\newcommand{\PSMathNotes}{}
\makeatletter
\NewDocumentCommand{\psplacednote}{m +m}{%
      \sbox{\PSNoteBox}{\parbox{\marginparwidth}{\footnotesize\raggedright\hand{\textsf{m#1.} #2}}}%
      \ifdim\dimexpr\ht\PSNoteBox+\dp\PSNoteBox\relax<\PSNoteLimit
        \leavevmode\marginpar{\usebox{\PSNoteBox}}%
      \else
        \par\nobreak\smallskip\hrule\nobreak\smallskip
        {\footnotesize\textsf{Marginal note m#1.}\quad\hand{#2}\par}
        \nobreak\smallskip\hrule\smallskip
      \fi}
\newcommand{\psflushmargins}{\ifmmode\else\ifinner\else\ifx\PSMathNotes\@empty\else\let\PSTempNotes\PSMathNotes\global\let\PSMathNotes\@empty\PSTempNotes\fi\fi\fi}
\NewDocumentCommand{\margin}{+m}{%
  \ifPSInApparatus\ifmmode\hand{\text{[marg.: }#1\text{]}}\else\hand{\textsf{[marg.: }#1\textsf{]}}\fi
  \else\stepcounter{psnote}%
    \ifmmode
      \text{\textsuperscript{m\arabic{psnote}}}%
      \begingroup\edef\PSQueue{\endgroup\noexpand\g@addto@macro\noexpand\PSMathNotes{\noexpand\psplacednote{\arabic{psnote}}{\unexpanded{$#1$}}}}\PSQueue
    \else
      \textsuperscript{m\arabic{psnote}}\psplacednote{\arabic{psnote}}{#1}%
    \fi
  \fi}
\AtBeginDocument{\everymath\expandafter{\the\everymath\aftergroup\psflushmargins}\everydisplay\expandafter{\the\everydisplay\aftergroup\psflushmargins}}
\makeatother
\newcommand{\psnumber}[2]{\expandafter\gdef\csname psnum@#1\endcsname{#2}}
\NewDocumentCommand{\var}{+m +m m m}{\ifnum\EditionMode=0 #1\else\ifnum\EditionMode=1\ifmmode\text{\textsuperscript{\tiny\csname psnum@#4\endcsname}}\else\textsuperscript{\tiny\csname psnum@#4\endcsname}\fi\fi #2\fi}
\newcommand{\psapparatusstart}{\par\bigskip\hrule\smallskip\PSInApparatustrue\begingroup\footnotesize\raggedright\textsf{Apparatus for scan \CurrentScan}\par}
\newcommand{\psapparatusend}{\par\endgroup\PSInApparatusfalse}
\NewDocumentCommand{\psentry}{m m m +m +m m}{\par\noindent\textbf{#2. #3.}\quad\textit{Typescript:} #4\quad\textit{Corrected:} \if\relax\detokenize{#5}\relax\textnormal{[omitted]}\else #5\fi\quad #6\par}
\newcommand{\psfrontmatter}{%
\hypersetup{pageanchor=false}\begin{titlepage}\raggedright
{\Large Alexander Grothendieck\par}\vspace{12mm}
{\Huge Pursuing Stacks\par}\medskip{\LARGE 1983\par}\vspace{12mm}
{\Large\EditionLabel\par}\vspace{8mm}
{\large\EditionStatus\par}\vfill
Université de Montpellier, Grothendieck archive, cotes 134-2--134-8.\par
Including the 26 source pages of R. Brown's translation at cote 134-2, sites 115--143.\par
\vspace*{2mm}\end{titlepage}\hypersetup{pageanchor=true}
\chapter*{Editorial statement}
AI-transcribed by Claude Opus 5.5 (Anthropic, effort high); checked against the scans, and this edition built, by GPT-6 Astra (OpenAI, ChatGPT Pro mode); prompts and records by Claude (Anthropic).

The witness consulted and not reproduced is the LaTeX edition of Mateo Carmona with the collaboration of Ulrik Buchholtz (arXiv 2111.01000 v2, 2021), extended from the 2015 edition by ``the scrivener'', CC0 1.0.

The translation in the appendix is R. Brown's, prepared with assistance from L. Breen and J-L. Loday, as his note says; the rights in it belong to its translator or his heirs.

The rights in the original work belong to its author or the author's representatives or heirs. Where we hold rights, in the AI-produced transcription, translation, corrections and typesetting, they are dedicated to the public domain under CC0 1.0. This edition is published for educational, non-profit purposes, and we claim no rights we do not hold. If you hold rights in this work and object to its publication here, please open a GitHub issue and we will take it down.

\section*{How to read this form}
\ifnum\EditionMode=0
The diplomatic form preserves typed and handwritten wording, cancellations, insertions, overwritten readings, underlining, doubts, and original page numbers. Blue marks handwriting; red labels handwritten cancellations; gray labels typewriter cancellations and underlying overwritten text. Angle brackets mark insertions; a superscript arrow separates an overwritten reading from its replacement. Typed insertions and typed overwriting remain plain. The disposition of long marginal notes is identified explicitly.
\else
The corrected text applies the author's insertions, replacements, cancellations, and recorded transpositions. It retains his English and notation, including consistent forms such as ``standart'', ``developped'', ``it's'' for ``its'', and ``A la''. An isolated slip is corrected only when specifically recorded and justified. Mathematical anomalies are preserved and discussed in the guide.
\ifnum\EditionMode=1
Blue continues to identify handwriting. A numbered call identifies each maximal editorial operation, with its full diplomatic and corrected readings and reason in the apparatus following the source page. Internal layers of a wholly replaced or deleted passage are preserved within that entry rather than duplicated as separate entries. These entries are generated from the same variants that govern the running text.
\else
This reading form omits the apparatus, cancelled material, superseded readings, insertion brackets, and color distinctions. It retains the resulting text, marginal notes, underlining, and unresolved doubts. Scan and source-page references appear in the margin.
\fi\fi

Every scan mark and remaining doubt links to its source page in the archive viewer. Doubts are editorial reading marks, not emendations. The source-page breaks identify the archival sequence; a source page may occupy more than one printed page. A partial working edition is explicitly labelled and contains only reviewed source pages.
\clearpage}

\begin{document}
\pssource{ps_134-7_081}{134-7/81}{https://grothendieck.umontpellier.fr/web/viewer.html?file=../134-7.pdf\#page=82}
\psnumber{ps_134-7_081-v0001}{1}
\psnumber{ps_134-7_081-v0002}{2}
\psnumber{ps_134-7_081-v0003}{3}
\psnumber{ps_134-7_081-v0004}{4}
\psnumber{ps_134-7_081-v0005}{5}
\psnumber{ps_134-7_081-v0006}{6}
\psnumber{ps_134-7_081-v0007}{7}
\psnumber{ps_134-7_081-v0008}{8}
\psnumber{ps_134-7_081-v0009}{9}
\psnumber{ps_134-7_081-v0010}{10}
\psnumber{ps_134-7_081-v0011}{11}
\psnumber{ps_134-7_081-v0012}{12}
\psnumber{ps_134-7_081-v0013}{13}
% Pursuing Stacks (cotes 134-2 to 134-8): archive cote 134-7, site page 81 (PDF page 82), typescript page 523.
% typescript. Draft compared: ps6.tex:4025-4111.
% AI-transcribed by Claude Opus 5.5 (claude-opus-5-5, adaptive thinking, effort high), Message Batches API, batch ps4 (msgbatch_012Aoc1NuskATPsyBeKSYK7J), 2026-11-12; reviewed against both supplied scan halves; see CHECK_DECISIONS.tsv.
\scanpage{134-7/81}
\mspage{523}
With this ih mind, we should'nt be too specific \var{\inserted{(for the time being)}}{\hand{(for the time being)}}{insertion}{ps_134-7_081-v0001} about the exact meaning
to be given to the sy\var{\xout{\illegible{1}}}{}{typewriter deletion}{ps_134-7_081-v0002}mbol\hand{s} Hotab$_0$(k) and Hot$_0$(k). It may be safer \var{\xout{too}}{}{typewriter deletion}{ps_134-7_081-v0003}
to replace these by Hotab$_1$(k) and Hot$_1$(k) (where the subscript 1 means
``1-connected''), all the more so as we hav'nt been able yet, in the
context of unipotent bundles, to give even a tentative precise definiti-%
on of Hot$_0$(k) i.e. of \hand{$\mathcal{M}$}$_0$(k) (mimicking the condition of abelian \var{\xout{f\illegible{1}}}{}{typewriter deletion}{ps_134-7_081-v0004} \hand{$\pi$}$_1$
with trivial action on the higher \hand{$\pi$}$_i$\hand{'}s), except by expressing things
via the associated ordinary homotopy type (passing ov er to X$_*$(k)), which
looks kind of stupid indeed ! But nothing is ``safe'' here anyhow, and
the subscript 0 looks definitely more natural here than subscript 1 ,
so we may as well keep it !

5.10.

What we are mainly interested in\hand{,} for expressing the relationships
between the non additive category Hot$_0$(k) and it's additive counterpart
Hotab$_0$(k), is an array of four functors \hand{\unsure{,}}in between these, two of
which being $L\overset{\hand{\sim}}{H}.$ (\var{\inserted{(total)}}{\hand{(total)}}{insertion}{ps_134-7_081-v0005}\var{\xout{t}}{}{typewriter deletion}{ps_134-7_081-v0006}``homology'',or ``\var{\deleted{abeliani}}{}{deletion}{ps_134-7_081-v0007}\allowbreak{}\var{\inserted{lineari}}{\hand{lineari}}{insertion}{ps_134-7_081-v0008}zation'') and $\overset{\hand{\sim}}{K}$ (Eilenberg-Mac-%
lane functor) in (1) and (2), the two remaining one's being L\hand{$\pi$}. (``total
homotopy'') and $\overset{\hand{\sim}}{S}$ (the ``spherical functor'' , compare p; 514 (9)), fit\var{\xout{tin}}{}{typewriter deletion}{ps_134-7_081-v0009}%
ting in\var{\inserted{to}}{\hand{to}}{insertion}{ps_134-7_081-v0010} the following diagram

\diagram{
% Hand-drawn arrows are marked blue; typed objects remain plain.
\begin{tikzcd}[column sep=huge]
\text{(4)}\qquad \text{Hotab}_0(k) \ar[color=blue,r, shift left, "\hand{\widetilde{K}}"] \ar[color=blue,r, shift right, "\hand{\widetilde{S}}"'] \& \text{Hot}_0(k) \ar[color=blue,l, bend right=40, "\hand{L\widetilde{H}.}"'] \ar[color=blue,l, bend left=40, "\hand{L\pi.}"]
\end{tikzcd}
}

I would now like to discuss the main formal properties to be expected
from this set of functors.

A) \uline{Adjunction} \uline{properties:}
\[
(5)\quad \hand{\left\{\rule{0pt}{2.5em}\right.}
\begin{array}{l}
\text{a) (LH.,}\hand{\widetilde{K}}\text{) is a pair of adjoint functors ,}\\
\text{b) (}\hand{\widetilde{S}}\text{, L\ .) is a pair of adjoint functors .}
\end{array}
\]

Thus we get adjunction homomorphisms (for L. in Hotab$_0$(k) and X
in Hot$_0$(k) )
\[
(6\ \text{a})\quad \hand{\left\{\rule{0pt}{2.5em}\right.}
\begin{array}{ll}
\text{X} \hand{\longrightarrow} \overset{\hand{\sim}}{\text{K}}(\text{L}\overset{\hand{\sim}}{\text{H}}.(\text{X})) & (\overset{\text{def}}{=} \overset{\hand{\sim}}{\text{W}}(\text{X}) \quad (\text{``pointed }\var{\deleted{\text{abeliani}}}{}{deletion}{ps_134-7_081-v0011}\var{\inserted{\text{lineari}}}{\hand{\text{lineari}}}{insertion}{ps_134-7_081-v0012}\text{zation''}))\\
\text{L}\overset{\hand{\sim}}{\text{H}}.(\overset{\hand{\sim}}{\text{K}}(\text{L.})) \hand{\longrightarrow} \text{L.} & ,
\end{array}
\]

and

\var{\xout{$S(L\illegible{1}(X))\illegible{1}X$}}{}{typewriter deletion}{ps_134-7_081-v0013}

\margin{81}

% Visual collation: both supplied image halves inspected; see CHECK_DECISIONS_134_7_8.tsv.

\psapparatusstart
\psentry{ps_134-7_081-v0001}{1}{insertion}{\inserted{(for the time being)}}{\hand{(for the time being)}}{The handwritten insertion is incorporated at its indicated position; insertion brackets leave the running text.}
\psentry{ps_134-7_081-v0002}{2}{typewriter deletion}{\xout{\illegible{1}}}{}{The typewriter strikes out this material; it leaves the corrected running text.}
\psentry{ps_134-7_081-v0003}{3}{typewriter deletion}{\xout{too}}{}{The typewriter strikes out this material; it leaves the corrected running text.}
\psentry{ps_134-7_081-v0004}{4}{typewriter deletion}{\xout{f\illegible{1}}}{}{The typewriter strikes out this material; it leaves the corrected running text.}
\psentry{ps_134-7_081-v0005}{5}{insertion}{\inserted{(total)}}{\hand{(total)}}{The handwritten insertion is incorporated at its indicated position; insertion brackets leave the running text.}
\psentry{ps_134-7_081-v0006}{6}{typewriter deletion}{\xout{t}}{}{The typewriter strikes out this material; it leaves the corrected running text.}
\psentry{ps_134-7_081-v0007}{7}{deletion}{\deleted{abeliani}}{}{The author cancels this material by hand; it leaves the corrected running text.}
\psentry{ps_134-7_081-v0008}{8}{insertion}{\inserted{lineari}}{\hand{lineari}}{The handwritten insertion is incorporated at its indicated position; insertion brackets leave the running text.}
\psentry{ps_134-7_081-v0009}{9}{typewriter deletion}{\xout{tin}}{}{The typewriter strikes out this material; it leaves the corrected running text.}
\psentry{ps_134-7_081-v0010}{10}{insertion}{\inserted{to}}{\hand{to}}{The handwritten insertion is incorporated at its indicated position; insertion brackets leave the running text.}
\psentry{ps_134-7_081-v0011}{11}{deletion}{$\deleted{\text{abeliani}}$}{}{The author cancels this material by hand; it leaves the corrected running text.}
\psentry{ps_134-7_081-v0012}{12}{insertion}{$\inserted{\text{lineari}}$}{$\hand{\text{lineari}}$}{The handwritten insertion is incorporated at its indicated position; insertion brackets leave the running text.}
\psentry{ps_134-7_081-v0013}{13}{typewriter deletion}{\xout{$S(L\illegible{1}(X))\illegible{1}X$}}{}{The typewriter strikes out this material; it leaves the corrected running text.}
\psapparatusend

\end{document}
